\section{Product Valuations and Kernel Composition}
\label{sec:monad}

The structure
through which classical probability supports sequential statistical
reasoning is the Giry monad:

\begin{definition}[Giry monad~{\cite{giry1982}}]
\label{def:giry}
The \emph{Giry monad} $G$ sends a measurable space $X$ to the space
$G(X)$ of probability measures on $X$.  Its unit sends a point
$x \in X$ to the point mass at $x$, and its multiplication sends a
measure $M$ on $G(X)$ to the measure
$A \mapsto \int m(A) \,\mathrm{d}M(m)$.
\end{definition}

\noindent
Probabilistic reasoning then composes in stages, and, when the monad
is commutative, independent stages compose in either order.  The
question is whether the intuitionistic calculus supports the same
structure.

A \emph{frame
homomorphism} is a map preserving finite meets and arbitrary joins.
Pushforward along
a frame homomorphism $f : \Omega \to \Omega'$ sends a valuation $v$
on $\Omega'$ to $a \mapsto v(f(a))$ on $\Omega$, and this preserves
the valuation axioms, making $\mathrm{Val}(-)$ a functor on
locales, the localic counterpart of the functor part of the Giry
monad~\cite{giry1982}.  A frame homomorphism is the opposite of a 
locale map, so the functor laws are checked against
identity and composition of frame
homomorphisms\leanref{valuation-comap}.

\begin{definition}[valued frames\leanref{valf}]
\label{def:spval}
The \emph{valuation functor} $\mathrm{Val} : \mathrm{Frm}^{\mathrm{op}} \to
\mathrm{Set}$ sends a frame to its set of valuations and a frame
homomorphism $\varphi : P \to Q$ to $w \mapsto w \circ \varphi$. The
category $\mathrm{SpVal}$ of \emph{valued frames} is its category of
elements: objects are pairs $(P, v)$ with $v$ a valuation on $P$, and a
morphism $(P, v) \to (Q, w)$ is a frame homomorphism $\varphi : P \to Q$
with $w \circ \varphi = v$.
\end{definition}

\noindent
A morphism of valued frames is a measure-preserving map of locales in
the opposite direction.  The point valuations
$\delta_p$ (Section~\ref{sec:representation}) provide the candidate
(monad) unit.
Theorem~\ref{thm:finite-rep} shows that, on finite frames,
$\mathrm{Val}$ agrees (on objects) with the classical
finite-distribution monad on points.  
Flattening a valuation of valuations is
integration, and integration in two variables requires product
valuations.  A \emph{kernel} is a map assigning a valuation to each point
of a locale, giving a conditional distribution.  By Kock's
identification of commutative monads with Fubini
theorems~\cite{kock2012}, kernels compose independently of order
exactly when $\mathrm{Val}$ is a commutative monad.  

A valuation $u$ on a product frame is a \emph{product} of
valuations $v$ and $w$ when $u(A \times B) = v(A) \cdot w(B)$ for
all rectangles $A \times B$.  Let
$\mathrm{Low}(\mathbb{N} \times \mathbb{N})$ be the frame of lower
sets of the grid $\mathbb{N} \times \mathbb{N}$, ordered
component-wise.  On this frame, products are unique:

\begin{theorem}[product uniqueness\leanref{valuation-product-unique}]
\label{thm:product-unique}
Let $v, w$ be valuations on $\mathrm{Low}(\mathbb{N})$.  Any two
valuations on $\mathrm{Low}(\mathbb{N} \times \mathbb{N})$ that are
products of $v$ and $w$ are equal.
\end{theorem}

\noindent
By the inclusion--exclusion equality, two valuations agreeing on a
$\sqcap$-closed family agree on all finite joins of its
members\leanref{eq-on-finsetsup-of-eq-on}.  The rectangles are
closed under $\sqcap$, and in
$\mathrm{Low}(\mathbb{N} \times \mathbb{N})$ every lower set is a
\emph{finite} union of
rectangles\leanref{exists-rect-decomposition}, so agreement on
rectangles is agreement everywhere.  Valuations carry no continuity
axiom (Section~\ref{sec:representation}), 
so the $\sigma$-additivity driving the
classical uniqueness argument for product measures has no analogue
here.  The order structure of the frame does that work instead.

\begin{theorem}[product non-uniqueness\leanref{fubini-fails-for-valuations}]
\label{thm:product-not-unique}
On the powerset frame $\mathrm{Set}(\mathbb{N} \times \mathbb{N})$
there exist valuations $v, w$ and two \emph{distinct} products of
$v$ and $w$.
\end{theorem}

\noindent
An \emph{ultrafilter} on $\mathbb{N}$ is a collection of subsets,
closed under intersection and superset, containing exactly one of
$S$ and its complement for every $S \subseteq \mathbb{N}$.  It is
\emph{free} when it contains no finite set.  The map assigning $1$
to the members of an ultrafilter and $0$ to every other subset is a
valuation\leanref{ultrafiltervaluation}.  A rectangle $S \times T$
belongs to an ultrafilter on the product exactly when $S$ and $T$
belong to its images under the two coordinate
projections\leanref{rect-mem-iff}.  A free ultrafilter
$\mathcal{U}$ on $\mathbb{N}$ yields two ultrafilters on
$\mathbb{N} \times \mathbb{N}$.  A set belongs to the former when,
for $\mathcal{U}$-most first coordinates, its slice contains
$\mathcal{U}$-most second coordinates.  A set belongs to the latter
under the same test with the coordinates in the opposite order.

The two ultrafilters' valuations agree on
\emph{every} rectangle yet differ on the upper triangle
$\{(m, m') \mid m < m'\}$\leanref{product-valuation-not-unique}.
Integrating the same marginal data in the two orders thus
disagrees.  Fubini fails for valuations, as it does for non-additive
integration~\cite{denneberg1994,ghirardato1997}.  By Kock's
correspondence, no canonical commutative
multiplication exists for unrestricted valuations.  This half of the
dichotomy is not specific to intuitionism, since powerset-frame
valuations are exactly the classical finitely additive probability
measures~\cite{bhaskara-rao1983}. Dropping continuity is what admits
the non-uniqueness. For continuous valuations on continuous domains the
product is unique and Fubini's theorem holds~\cite{jones1990}, and they
form the probabilistic powerdomain~\cite{jones-plotkin1989}.  A monad therefore requires restricting either
the frames or the valuations.

A first step in the second direction is to restrict the mixing
measures to countable ones. The components of a mixture are arbitrary
valuations, so this restricts how valuations are combined, not which
valuations occur. With countable weights the order of mixing does not
matter, by the interchange of double series. In
Theorem~\ref{thm:product-not-unique} the mixing measures are ultrafilter
valuations, which are not countable mixtures:

\begin{proposition}[countable kernel composition\leanref{mixcountable}]
\label{prop:mixcountable}
\begin{enumerate}
\item A countable mixture of valuations is a
  valuation\leanref{mixcountable}.
\item A mixture of mixtures flattens to a mixture with product
  weights, which is discrete
  Chapman--Kolmogorov\leanref{mixcountable-mixcountable}.
\item The one-point mixture is the
  identity\leanref{mixcountable-unique}.
\item Iterated countable mixing is
  order-independent\leanref{mixcountable-comm}.
\end{enumerate}
\end{proposition}

\noindent
Restricted to countable-mixture presentations, the data of a monad are
all in place: the category of locales, $\mathrm{Val}$ restricted to
countable mixtures as the functor, the point valuations as the unit,
and countable mixing as the multiplication.
Proposition~\ref{prop:mixcountable} verifies exactly the laws this
needs: item~1 keeps the multiplication inside the fragment, item~2 is
its associativity, item~3 its unit law, and item~4 its commutativity.
What is not yet settled is whether this descends to valuations
themselves, since a valuation can admit more than one countable-mixture
presentation, and the diffuse part of Section~\ref{sec:representation}
is not captured by any countable presentation at all.
